% Macros for content
\newcommand{\systemname}{Your System}
\newcommand{\teamname}{Your Team}
\newcommand{\workshopdate}{\today}
\newcommand{\iteration}{v1.0}

\newcommand{\businesscaseheadline}{Business Case}
\newcommand{\businesscasecontent}{
Brief description of the business case or economic driver behind the software system.
}
\newcommand{\businesscasehint}{
If the business needs to make money, there is an economic driver and a business case behind every software initiative. If not, it’s a hobby or an NPO project.

And it all start with WHY. - Start with WHY.

A common pitfall is that the development team, or worse, the entire product team, is unaware of the specific economic drivers and business case behind these software initiatives.
}

\newcommand{\functionaloverviewheadline}{Functional Overview}
\newcommand{\functionaloverviewcontent}{
The most important functional requirements at a high level.
}
\newcommand{\functionaloverviewhint}{
As you work out the functional overview of the software initiative as a product team, imagine a product box of your software - what is on your software box?

What is on your software product box?

Think the product box to get the answers to the following questions:

\begin{itemize}
\item What business value does the software deliver?
\item What are the higher-level business functions?
\end{itemize}

With the answers to these questions, you can describe the key functional requirements at a high level.
}

\newcommand{\qualitygoalsheadline}{Quality Goals}
\newcommand{\qualitygoalscontent}{
The three most important quality goals for the architecture, which have the highest priority for the most important stakeholder.
}
\newcommand{\qualitygoalshint}{
Here you work out the three most important quality goals of the stakeholders for the software system.

It’s important that these quality goals are known by every team member because they shape your architecture.

The standard \href{https://iso25000.com/index.php/en/iso-25000-standards/iso-25010}{ISO/IEC 25010} provides an overview of possible quality attributes of interest, the modern arc42 quality model (\href{https://quality.arc42.org/}{Q42}) collects more than 130 quality attributes with extensive explanation and examples.
}

\newcommand{\businesscontextheadline}{Business Context}
\newcommand{\businesscontextcontent}{
Separate your system under construction as a black box from all its communication partners. Communication partners are neighbouring external systems and users.
}
\newcommand{\businesscontexthint}{
Business context focuses on your system as a black box. Attention is focused on the system’s environment and how the system interacts with all external systems within it.

With the business context, you can identify the communication paths and thus potential technical risks.

It’s also a good way to discuss with various stakeholders the scope of the system and the communication paths to various external systems.

When trying to work out the business context, try to identify the users of your system and all neighbouring systems.
}

\newcommand{\orgconstraintsheadline}{Organisational Constraints}
\newcommand{\orgconstraintscontent}{
Any organisational requirement that limits the software architects' freedom of decision.
}
\newcommand{\orgconstraintshint}{
Constraints determine the degree of your freedom to make decisions as an engineering team.

In this part you work out all the requirements that limit your freedom as a product team in developing the (software) product.

Organisational constraints often include time and money. We’ve also seen that there are constraints on the methodologies used (such as SAFe, Scrum, …) or techniques.
}

\newcommand{\techconstraintsheadline}{Technical Constraints}
\newcommand{\techconstraintscontent}{
Any technical requirement that limits the software architects' freedom of decision.
}
\newcommand{\techconstraintshint}{
Constraints determine the degree of your freedom to make decisions as an engineering team.

In this part you work out all the requirements that limit your freedom as a product team in developing the (software) product.

In larger enterprises, there are often a lot of technical constraints. There is often an overarching architecture committee that sets the standards for technology choices.

For example, if you decide to build a single-page application, the overarching architecture committee will set Angular as the SPA framework, so there are expertise synergies between teams.
}

\newcommand{\archhypothesesheadline}{Architectural Hypotheses}
\newcommand{\archhypothesescontent}{
Resulting architectural hypotheses and important, expensive, large-scale or risky decisions, including justifications.
}
\newcommand{\archhypotheseshint}{
Congratulations - you’ve created your architectural playground.

Now you can create your first architectural hypotheses based on your current team knowledge.

Be aware: This is about higher-level architectural assumptions, not micro-decisions like design questions.

Ideally, you can turn your architectural assumptions into draft architecture decision records.
}

\newcommand{\challengesrisksheadline}{Technical Challenges \& Risks}
\newcommand{\challengesriskscontent}{
Identified current known challenges and technical risks.
}
\newcommand{\challengesriskshint}{
Based on the playground you have created, such as business context and scope, quality goals and constraints, you can capture the upcoming challenges and the resulting risks.

It’s very important to work this out, evaluate it continuously, and communicate it to the key stakeholders in the system.

This section answers the question of how you evaluate the current situation in relation to the elaborated architectural playground.
}
